% ============================================================
%  ANSWER KEY — numbered as in chapters/exercises.tex
% ============================================================
\clearpage
\section{Answer Key}

Male-speaker forms are given; a female speaker uses the \textit{-ii} forms. In the future and subjunctive, \hw{ँ} and \hw{ं} are both acceptable spellings (\hw{जाऊँगा} = \hw{जाऊंगा}).

\subsection{4.1 Present habitual \& continuous}
1.\ \hw{पीता} \quad 2.\ \hw{रहती} \quad 3.\ \hw{सीख रहे} \quad 4.\ \hw{कर रहे} \quad 5.\ \hw{खेलते} \quad 6.\ \hw{बना रही}

\subsection{4.2 Simple past}
\begin{enumerate}
  \item \hw{मैं ऑफिस गया।} — \textit{jaanaa} is intransitive: no \textit{ne}.
  \item \hw{मैंने चाय पी।} — agrees with \hw{चाय} (f.).
  \item \hw{वह घर आई।} — intransitive; agrees with the subject (f.).
  \item \hw{हमने खाना खाया।} — agrees with \hw{खाना} (m.sg.).
  \item \hw{मैंने किताब पढ़ी।} — agrees with \hw{किताब} (f.).
  \item \hw{मैंने साइकिल चलाई।} — agrees with \hw{साइकिल} (f.).
  \item \hw{उसने काम किया।} — \hw{वह} + \textit{ne} = \hw{उसने}; irregular \hw{किया}.
\end{enumerate}

\subsection{4.3 Future tense}
1.\ \hw{जाऊँगा} \quad 2.\ \hw{बनाएगी} \quad 3.\ \hw{आओगे} \quad 4.\ \hw{सीखेंगे} \quad 5.\ \hw{पिएँगे} \quad 6.\ \hw{रहेगा}

\subsection{4.4 Oblique case}
1.\ \hw{कमरे} \quad 2.\ \hw{लड़के} \quad 3.\ \hw{दोस्तों} \quad 4.\ \hw{बड़े शहर} \quad 5.\ \hw{किताबों} \quad 6.\ \hw{लड़कियों}

\subsection{4.5 Possessives \& \textit{apnaa}}
1.\ \hw{अपने} \quad 2.\ \hw{मेरा} \quad 3.\ \hw{अपनी} \quad 4.\ \hw{आपकी} \quad 5.\ \hw{मेरी} \quad 6.\ \hw{अपने}

\noindent In 2 and 5 the owner (\textit{main}) is not the subject (\hw{नाम}, \hw{बहन} are), so \textit{apnaa} is not used.

\subsection{4.6 Imperative}
1.\ \hw{बैठिए / बैठो} \quad 2.\ \hw{आइए / आओ} \quad 3.\ \hw{कीजिए / करो} \quad 4.\ \hw{दीजिए / दो} \quad 5.\ \hw{बोलिए / बोलो} \quad 6.\ \hw{पीजिए / पियो} \quad 7.\ \hw{मत जाओ!}

\subsection{4.7 Joining actions with \textit{-kar}}
\begin{enumerate}
  \item \hw{मैं नाश्ता करके ऑफिस जाता हूँ।} — \textit{karnaa} → \textit{karke}.
  \item \hw{हाथ धोकर खाइए।}
  \item \hw{मैं खाना खाकर घर गया।} — the last verb (\textit{jaanaa}) is intransitive, so \hw{मैं}, not \hw{मैंने}.
  \item \hw{उसने घर आकर चाय पी।} — the last verb (\textit{piinaa}) is transitive, so \hw{उसने}.
\end{enumerate}

\subsection{4.8 Subjunctive \& conditionals}
\begin{enumerate}
  \item \hw{क्या मैं अंदर आऊँ?}
  \item \hw{चलें!} (colloquial: \hw{चलो!})
  \item \hw{मैं क्या करूँ?}
  \item \hw{शायद वह कल आए।}
  \item \hw{अगर बारिश होगी, तो मैं घर पर रहूँगा।}
  \item \hw{अगर आप थके हैं, तो आराम कीजिए।}
\end{enumerate}

\subsection{5.1 Masculine or feminine?}
1.\ f. \quad 2.\ m. \quad 3.\ f. \quad 4.\ m. \quad 5.\ f. \quad 6.\ m. \quad 7.\ f. \quad 8.\ m. \quad 9.\ f. \quad 10.\ m. \quad 11.\ f. \quad 12.\ m.

\subsection{5.2 Write the Hindi word}
1.\ \hw{सुबह} \quad 2.\ \hw{मंगलवार} \quad 3.\ \hw{दादी} \quad 4.\ \hw{महँगा} \quad 5.\ \hw{बाएँ} \quad 6.\ \hw{लाल} \quad 7.\ \hw{मौसम} \quad 8.\ \hw{दोस्त} \quad 9.\ \hw{कल} \quad 10.\ \hw{सिर} \quad 11.\ \hw{साइकिल} \quad 12.\ \hw{फूल}

\subsection{5.3 Opposites}
1.\ \hw{छोटा} \quad 2.\ \hw{बुरा} \quad 3.\ \hw{ठंडा} \quad 4.\ \hw{सस्ता} \quad 5.\ \hw{पुराना} \quad 6.\ \hw{बाएँ} \quad 7.\ \hw{नीचे} \quad 8.\ \hw{दूर}

\subsection{5.4 Numbers}
1.\ \hw{सात} \quad 2.\ \hw{बारह} \quad 3.\ \hw{पंद्रह} \quad 4.\ \hw{बीस} \quad 5.\ \hw{पच्चीस} \quad 6.\ \hw{तीस} \quad 7.\ \hw{पचास} \quad 8.\ \hw{सौ}

\subsection{5.5 What time is it?}
\begin{enumerate}
  \item \hw{तीन बजे हैं।}
  \item \hw{एक बजा है।} — singular for one o'clock.
  \item \hw{साढ़े चार बजे हैं।}
  \item \hw{डेढ़ बजा है।}
  \item \hw{ढाई बजे हैं।}
  \item \hw{सवा छह बजे हैं।}
  \item \hw{पौने आठ बजे हैं।}
  \item \hw{पाँच बजकर दस मिनट हुए हैं।}
\end{enumerate}

\subsection{5.6 Reading practice}
1.\ namaste \quad 2.\ kitaab \quad 3.\ dost \quad 4.\ laRkii \quad 5.\ hindii \quad 6.\ kaun \quad 7.\ ThanDaa \quad 8.\ kripayaa \quad 9.\ sTeshan \quad 10.\ mausam \quad 11.\ pahaaR \quad 12.\ zyaadaa

\subsection{5.7 Reading a text message}
\textit{Translation:} Can we meet this Sunday? Tomorrow I'm taking part in a 30 km race, and I'm still on the way to Dharamshala. Thank you.
\begin{enumerate}
  \item This Sunday (\hw{इस रविवार को}).
  \item She is running a 30 km race. \emph{kal} = tomorrow, because the verb is present continuous (a plan), not past.
  \item A woman: \hw{रही} (a man would write \hw{रहा}).
  \item On the way to Dharamshala (\hw{धर्मशाला के रास्ते में}).
  \item \hw{दौड़} is feminine → \hw{की}; \hw{रास्ता} is masculine and oblique before \hw{में} → \hw{के रास्ते}.
  \item \hw{मैं कल 30 किलोमीटर की दौड़ में भाग ले रहा हूँ।} (only \hw{रही} → \hw{रहा} changes)
  \item \hw{हाँ, हम रविवार को मिल सकते हैं। शुभकामनाएँ!}
\end{enumerate}

\subsection{6.1 Find the mistake}
\begin{enumerate}
  \item \hw{मैंने साइकिल चलाई।} — transitive past needs \textit{ne}.
  \item \hw{मैं अपने घर जाता हूँ।} — owner = subject → \textit{apnaa}.
  \item \hw{आप क्या कर रहे हैं?} — \textit{aap} takes \hw{हैं}.
  \item \hw{वह तैर नहीं सकती।} — \hw{नहीं} goes before \textit{saknaa}.
  \item \hw{मैं ऑफिस गया।} — \textit{jaanaa} is intransitive: no \textit{ne}.
  \item \hw{लड़के को पानी दो।} — oblique before \hw{को}.
  \item \hw{मुझे चाय पसंद है।} — \textit{pasand} agrees with the thing liked, not with \textit{main}.
  \item \hw{मैंने कुट्टू के फूल देखे हैं।} — agrees with \hw{फूल} (m.pl.).
\end{enumerate}

\subsection{6.2 English → Hindi}
\begin{enumerate}
  \item \hw{मैं हिंदी सीख रहा हूँ।}
  \item \hw{मेरे पास साइकिल है।}
  \item \hw{मुझे पानी चाहिए।}
  \item \hw{धीरे बोलिए।} (or \hw{कृपया धीरे बोलिए।})
  \item \hw{कल मैं बस से ऑफिस गया।}
  \item \hw{मैंने कभी कुट्टू नहीं खाया है।}
  \item \hw{मुझे कल काम करना होगा।}
  \item \hw{मेरी बहन मुझसे बड़ी है।}
  \item \hw{अगर तुम आओगे, तो हम साथ खाएँगे।}
  \item \hw{मेरे सिर में दर्द है।}
  \item \hw{क्लास शाम छह बजे है।}
  \item \hw{आज मैंने शक्ति से अघम तक साइकिल चलाई।}
\end{enumerate}

\subsection{6.3 Hindi → English}
\begin{enumerate}
  \item I like mangoes.
  \item How far is the station from here?
  \item I didn't have breakfast.
  \item She used to live in a village.
  \item I ate and then went home. (lit.\ having eaten, I went home)
  \item May I come in? / Shall I come in?
  \item It's very cold today.
  \item Please say it again.
\end{enumerate}
